\chapter{Conclusions}
\label{ch:conclusions}

This thesis has documented the complete engineering cycle of a modular robotic limb, moving from the theoretical
kinematic analysis to the physical realization of a high-performance, cost-effective actuation system. The central
challenge addressed in this work was the "actuation bottleneck" typical of low-cost robotics, where the mechanical
and control performance is severely limited by the poor quality of commercial hobbyist servo-motors. By adopting a
top-down design methodology, supported by simulation and component-level redesign, this research successfully
demonstrated that high-fidelity control can be achieved without resorting to prohibitively expensive industrial
hardware.


\section{Achievements and Considerations}

The primary technical achievement of this work is the development and validation of the custom "Smart Actuator" module.
The transition from a standard analog architecture to a fully digital, microcontroller-based system yielded quantifiable
improvements in every key performance metric.

As demonstrated by the experimental benchmarking in Chapter \ref{ch:servomotor}, the custom module reduced the Root
Mean Square Error (RMSE) in position tracking from an average of $\approx 14^{\circ}$ (observed in the original MG996R
units) to just $0.37^{\circ}$. This result confirms that the limitations of low-cost robotics are often not mechanical
but electronic and algorithmic. Specifically:

\begin{itemize}
    \item \textbf{Linearity and Precision:} The integration of the AS5600 magnetic encoder eliminated the
        non-linearities and wear associated with resistive potentiometers, providing a stable 12-bit feedback
        signal.
    \item \textbf{Friction Compensation:} The implementation of the Switching PID control strategy proved effective in
        mitigating the stick-slip phenomenon. The ability to dynamically switch integrator states allowed the actuator
        to overcome static friction without inducing the limit-cycle oscillations (hunting) typical of standard linear
        controllers.
    \item \textbf{Sim-to-Real Fidelity:} The Hardware-in-the-Loop validation detailed in Chapter \ref{ch:prototype}
        confirmed a high degree of correlation between the Simulink numerical model and the physical prototype.
\end{itemize}


\section{Economical and Practical Aspects}

From an economical perspective, the custom module represents a strategic compromise between cost and performance.
Industrial actuators capable of similar precision typically cost between 10 to 20 times more than a standard hobby
servo. The prototype developed in this thesis achieves comparable tracking fidelity at a fraction of that cost,
maintaining a bill of materials (BOM, Appendix \ref{app:bom}) compatible with academic and hobbyist budgets.

A crucial practical advantage of this design is its retro-compatibility. By strictly adhering to the external
form factor of the standard "40x20x36mm" servo case, the custom module acts as a "drop-in replacement". This allows for
the immediate upgrade of existing hexapod platforms without requiring any mechanical redesign or the re-printing of
structural parts. Furthermore, the shift to an $I^2C$ bus architecture simplifies the robot's wiring harness, allowing
multiple actuators to be daisy-chained on a single cable, thereby reducing the complexity and weight of the central
electronics.


\section{General Applicability and Future Work}

While this research focused on a hexapod limb, the developed "Smart Actuator" has general applicability across the field
of robotics. It can be effectively employed in low-cost robotic arms, camera gimbals, or bipedal systems where precision
and customizability are required but budget is constrained.

However, several areas remain open for future development to transition from this single-leg prototype to a fully
autonomous hexapod robot:

\begin{itemize}
    \item \textbf{Proprioceptive Validation:} The adaptive "Blind Walking" strategy and the Finite State Machine
        trajectory generator were validated only in the simulation environment (Chapter \ref{ch:model}). Future work
        must focus on enabling the impact detection logic on the physical hardware, utilizing the motor current feedback
        or external contact sensors to handle uneven terrain.
    \item \textbf{Dynamic Tuning Optimization:} The PID gains used in the experimental validation were tuned empirically.
        Implementing an auto-tuning routine or a model-based parameter identification algorithm would further optimize
        the controller's response to varying loads.
    \item \textbf{Full-Body Integration:} The final step is the replication of the limb design to assemble the complete
        hexapod. This will introduce new challenges regarding power management, inter-leg coordination, and the global
        stability of the chassis, which were outside the scope of this single-leg study.
\end{itemize}

In conclusion, this thesis lays a solid foundation for the development of accessible, high-performance legged robots,
proving that intelligent control design can effectively bridge the gap between low-cost hardware and advanced robotic
capabilities.